In my weld-seam HMI, I used one modified-DH transformation chain for analytical kinematics, joint controls, and an OpenGL digital twin. During CR5 integration, I documented the accuracy limits of a usable eye-to-hand transform for fixed-view data collection and inference, checked positions and targets against workspace bounds before using or replaying data, and adopted blocking short-horizon ServoP to preserve state-action consistency. I now want to connect these decisions within a stronger framework for precision, sensing, control, and mechatronics. HKUST's MSc in Mechanical Engineering would provide that framework.

If the current courses run in my intake and MESF 6950 Independent Project is approved, I would organize these resources around an instrumented robot prototype. MESF 5550 Robotics would deepen the robot and control reasoning behind it; MESF 5560 Precision Manufacturing Technologies would clarify how mechanical precision sets meaningful operating tolerances; and MESF 5970 Smart Materials, Sensors and Transducers would broaden its sensing and actuation basis. MESF 6950 could then bring geometry, calibration, sensing, workspace limits, and execution tests into the same system. My experience connecting coordinate models to observed motion would provide a concrete starting point, while measured responses would reveal which physical assumption needs revision. This preparation would support robotics or intelligent-manufacturing R\&D in which mechanical design, sensing, and control develop together.
